\section{Evaluation Policy Register}
\label{app:policy-register}

This appendix records the frozen policy dimensions. The versioned release
document names a base profile and may override only schema-authorized top-level
bounds and blocked extensions. Any other edit changes the implementation commit
and creates a new run.

\subsection{Registered profiles and evaluation selection}
\label{subsec:registered-profiles}

The implementation defines staging, strict, and paranoid profiles, with public,
enterprise, and high-risk aliases. Staging can warn or quarantine and is only a
rehearsal profile. Strict is the primary release base. It requires supplied
type evidence, blocks unknown kinds and decoder failures, enforces output
convergence, scans input and output, and requires semantic reconstruction for
static images, GIF, audio, and video. Paranoid narrows budgets and selected
output sets. A paranoid comparison is paired and not pooled with strict. Dry
run, warning, bypass, and unmodified fallback paths never count as accepted.

\subsection{Security-policy fields}
\label{subsec:security-policy-fields}

\Cref{tab:policy-register} identifies whether each field is source-defined or a
schema-authorized release-document value and states its evidence obligation.

\begingroup
\footnotesize
\setlength{\tabcolsep}{5pt}
\begin{longtable}{L{0.20\textwidth}L{0.14\textwidth}L{0.23\textwidth}L{0.24\textwidth}}
\caption{Security-policy register and required evidence.}
\label{tab:policy-register}\\
\toprule
Field & Owner & Release rule & Required evidence \\
\midrule
\endfirsthead
\toprule
Field & Owner & Release rule & Required evidence \\
\midrule
\endhead
Profile and version & Release document & Known version and strict base & Parsed document, schema validation, SHA-256 \\
Input bytes & Release document & Reject above bound before handler work & Cases below, at, and above $L$, with stage record \\
Output bytes & Release document & Terminate or reject above bound & Child-output and in-process boundary cases \\
Expansion ratio & Release document & Reject candidate ratio above bound & Constructible below, at, and above cases \\
Handler duration & Source & Earliest of policy and caller deadline & Expired, near-bound, timeout, and cancellation cases \\
Required MIME & Source & Present and recognized in strict profile & Missing, alias, unsupported, and conflicting cases \\
Type agreement & Source & Name, supplied type, bytes, output type, and parser converge & S3 crossing matrix and T3 report \\
Unknown kind & Source & Block & Unknown and unsupported S6 records \\
Decoder or probe failure & Source & Block in strict profile & Missing tool, invalid decode, timeout, and nonzero exit \\
Reconstruction level & Source by kind & Structural or semantic-or-block as registered & Per-record performed and required fields \\
Signature policy & Source and rule input & Input and output scan, blocking pre-scan in strict & Rule digest, location, severity, and stage result \\
Blocked extensions & Source plus release document & Normalized union, no silent deletion & Name vectors and serialized normalized policy \\
Output MIME allowlist & Source & Enforce per file kind & Allowed, denied, unknown, and cross-kind outputs \\
Image pixels & Source & Check before and after decode & Dimension and aggregate-pixel boundaries \\
Animation frames and pixels & Source & Bound frames, frame area, aggregate work, and duration & S5 generated boundary records \\
Audio duration, bitrate, channels & Source & Probe then enforce before semantic processing & Independent probe record and boundary cases \\
Video duration and bitrate & Source & Probe then enforce before semantic processing & Independent probe record and boundary cases \\
Container elements and nesting & Source validator & Exact grammar-specific budgets & S2 and S5 mutation records \\
External process protocols & Source & Local file and pipe protocols only & Command capture and prohibited-protocol regression \\
External process diagnostics & Source & Bounded and sanitized & Oversized output and control-character injection \\
Publication & Source & Atomic publication for Clean only & Fault campaign directory trace and hashes \\
\bottomrule
\end{longtable}
\endgroup

Bounds use implementation units. Ratios derive from exact integer sizes. The
serialized value, not a manuscript approximation, is authoritative.

\subsection{Reconstruction register}
\label{subsec:reconstruction-register}

The required level is checked against the handler report. Structural processing
supports container and metadata claims but not neutralization inside retained
compressed payloads. Semantic processing creates output from bounded decoded
state and remains conditional on codec assumptions. Semantic-or-block forbids
structural fallback.

\begin{table}[tb]
\centering
\caption{Evaluation reconstruction requirements and admissible claims.}
\label{tab:reconstruction-register}
\small
\begin{tabularx}{\textwidth}{L{0.19\textwidth}L{0.20\textwidth}Y L{0.25\textwidth}}
\toprule
Family & Strict requirement & Admissible accepted-output claim & Explicit exclusion \\
\midrule
Static raster images & Semantic or block & New canonical image from bounded decoded pixels & Harmful visual meaning and steganographic content \\
GIF & Semantic or block & New bounded frame sequence with documented timing normalization & Arbitrary meaning encoded in frames \\
APNG & Semantic or block for high-risk claim & New semantic animation when available & Structural remux is insufficient for bitstream disarm \\
Animated WebP & Semantic or block for high-risk claim & New semantic animation when available & Structural RIFF reconstruction and complete frame decoding do not replace retained compressed frame payloads \\
MP3, WAV, FLAC, Ogg & Semantic or block in strict audio profile & New output from decoded audio representation & Structural copy cannot neutralize retained codec payload \\
MP4, MOV, M4A & Semantic or block in strict audio or video profile & Constrained semantic output after external decoding & Structural atom remux supports only container claims \\
WebM & Semantic or block in strict video profile & Constrained semantic output after external decoding & Structural EBML remux supports only container claims \\
Unsupported classes & Not applicable & No deliverable under media-only policy & No behavioral family classification \\
\bottomrule
\end{tabularx}
\end{table}

\subsection{Policy ordering for fixed-witness regression}
\label{subsec:policy-ordering}

The relation $p_2\sqsubseteq p_1$ means that $p_2$ is at least as restrictive
as $p_1$ under one implementation and parser-registry version. It requires no
larger maximum bounds, no smaller minimum requirements, no weaker
reconstruction, no broader output language, no smaller rejection relation, and
no bypass in $p_2$. The regression holds candidate bytes and parser reports
fixed. A policy that produces another candidate is reported as a reconstruction
change rather than a monotonicity result.

Paranoid $\sqsubseteq$ strict applies only after a generated policy diff. Staging is
unordered because warnings and quarantine change terminal semantics.

\subsection{Independent fidelity policy}
\label{subsec:fidelity-policy-register}

Fidelity thresholds are held in a separate schema-valid document. That document
maps each output profile to a metric profile, required tools, required structural
comparisons, objective quality metrics, and pass thresholds. It also identifies
documented transformations such as metadata removal, codec conversion, fixed
application identifiers, or animation delay normalization. The document is
frozen before corpus execution and its SHA-256 is included in the fidelity
report.

Missing metrics, nonfinite values, decoder errors, and unevaluated accepted
outputs fail fidelity. Thresholds are justified and frozen before results. No
observed metric is asserted here.

\subsection{Change control}
\label{subsec:policy-change-control}

Digests cover release, alias, fidelity, format, and tool configuration. Any
post-run change invalidates the old interpretation and requires a new identifier
and complete rerun of the controlled sequence in
\cref{fig:evaluation-pipeline}.
